% template-only:begin
% Optional background chapter: a literature review by theme, the theoretical framework, and
% data or methods shared by several research chapters. It is not part of the thesis by default.
% To use it, copy this file by hand to chapters/background.tex, delete this block, and
% uncomment `% \include{chapters/background}` in main.tex (below the introduction). Its labels
% are fixed (ch:background, sec:background-*) because a thesis has at most one such chapter.
% Delete the sections the thesis does not need.
% template-only:end
\chapter{Background}
\label{ch:background}

% With this chapter, the general introduction keeps only a short overview of the field and
% the gap, and refers here (\cref{ch:background}) for the full review.

\section{Literature Review}
\label{sec:background-literature}
% TODO: Review the literature by theme or question, not paper by paper. For each theme say
% what is established, what is contested, and what is missing. Examiners look for a critical
% synthesis that covers the key and the recent work of the field.
\mbox{}\par

\section{Theoretical Framework}
\label{sec:background-framework}
% TODO: The concepts, models or theory that the studies rely on, with the definitions and
% notation used throughout the thesis (keep notation.tex and the list of symbols in step).
\mbox{}\par

\section{Shared Data and Methods}
\label{sec:background-methods}
% TODO: Data sets, instruments and methods used by more than one research chapter, described
% once here; each research chapter refers back with \cref{sec:background-methods} and describes
% only what is specific to it.
\mbox{}\par

\section{Research Gaps and Research Questions}
\label{sec:background-gaps}
% TODO: Summarise the gaps that the review exposed and how they lead to the research
% questions (the same RQ ids as \cref{sec:introduction-aims} and docs/thesis-plan.md).
\mbox{}\par
